\begin{center}
\textsc{Gauss} to \textsc{Bessel}.\quad Göttingen, 1~April 1827.
\end{center}

\par\noindent\hspace*{\parindent}.\ .\ .\ .\ .\ .\quad This spring I intend to measure the amplitude of the arc between
Göttingen and Altona with the zenith sector, and I myself am eager
to see the result. A series of observations made with the meridian circles at the beginning of
1824 gave \(2^{0}\,1^{\prime}\,58^{\prime\prime}\). The altitudes of the pole, taking mine as I adopt it from the
best observations and \textsc{Schumacher}'s as he gives it, yield
more like 1 or \(1\frac{1}{2}\) arcseconds less. The geodetic measurements, however,
assuming the uniform shape of the Earth and \textsc{Walbeck}'s dimen-
\par\clearpage
\noindent{\small\textbf{ON THE HANOVER TRIANGULATION.}}\hfill363\par\medskip
\noindent sions, give 4 or 5 arcseconds more. Almost exactly the same value as this
last result is given by the exactly simultaneous zenith distances of zenith stars
that \textsc{Nehus} in Altona and I here observed with the meridian circle, with
the zero point determined by collimators. \textsc{Nehus} is now here to help me
with the sector observations. He has brought a \textsc{Repsold} suspended colli\-%
mator, but so far I consider the method of finding the nadir point by
mercury reflection considerably more accurate. I think, however, that the former method
can still be improved.\quad.\ .\ .\ .\ .\ .
